\documentclass[11pt]{article}
\usepackage[utf8]{inputenc}
\usepackage{amsmath,amssymb,amsthm}
\usepackage{hyperref}
\usepackage{graphicx}
\usepackage{booktabs}
\usepackage{listings}
\usepackage{array}
\usepackage[margin=1in]{geometry}

\newcommand{\R}{\mathbb{R}}

\lstset{
  basicstyle=\ttfamily\small, breaklines=true, frame=single,
  extendedchars=true,
  literate={ô}{{\^o}}1 {ℝ}{$\mathbb{R}$}1 {β}{$\beta$}1 {σ}{$\sigma$}1 {θ}{$\theta$}1
           {α}{$\alpha$}1 {⟶}{$\to$}1 {→}{$\to$}1 {⊗}{$\otimes$}1 {≤}{$\leq$}1
           {≥}{$\geq$}1 {≠}{$\neq$}1 {≃}{$\simeq$}1 {⟨}{$\langle$}1 {⟩}{$\rangle$}1
           {₀}{$_0$}1 {₁}{$_1$}1 {₂}{$_2$}1 {₃}{$_3$}1 {₄}{$_4$}1
           {ⁿ}{$^n$}1 {ε}{$\varepsilon$}1 {μ}{$\mu$}1 {ν}{$\nu$}1
           {λ}{$\lambda$}1 {Λ}{$\Lambda$}1 {Δ}{$\Delta$}1 {⊢}{$\vdash$}1
}

%% Bundle target: I3 (tier 3, journal JOSS | CPC)
%% Schema: docs/BUNDLE_DIRECTORY_SCHEMA.md
%% Source manifest: papers/I3/source_manifest.md (sourceless from
%%   PAPER_DRAFT_MAPPING.md per Phase 7a sub-wave 7a.3 + Phase 6o.zeta)
%% Append log: papers/I3/append_log.json
%% Change log: papers/I3/change_log.md
%% This bundle is fresh-authored from Lean-module substrate (Phase 6o
%%   Wave 3b: 12 Lean modules under Itô/* + LDP/*; cross-bridge to
%%   Phase 6n Wave 2a Glorioso-Liu axiomatic + Phase 6n Wave 2c.5c+
%%   IsLDPRateFunction typeclass).
% Auto-generated by update_counts.py — 2026-04-24T19:44:03.140400
% Do not edit manually. Run: uv run python scripts/update_counts.py
\newcommand{\totaltheorems}{3993}
\newcommand{\substantivetheorems}{3882}
\newcommand{\placeholdertheorems}{111}
\newcommand{\axiomcount}{1}
\newcommand{\sorrycount}{0}
\newcommand{\sorrypercent}{0.0}
\newcommand{\leanmodules}{167}
\newcommand{\leandefinitions}{3297}
\newcommand{\pythonmodules}{68}
\newcommand{\testfiles}{59}
\newcommand{\totaltests}{2836}
\newcommand{\figurecount}{111}
\newcommand{\notebookcount}{54}
\newcommand{\papercount}{19}
\newcommand{\aristotleproved}{322}
\newcommand{\aristotleruns}{44}
% --- Paper 18 / Phase 5t per-section counts (FermiHubbardDimer.lean) ---
\newcommand{\fhdTotal}{143}
\newcommand{\fhdWaveTwo}{13}
\newcommand{\fhdWaveThree}{6}
\newcommand{\fhdWaveFour}{22}
\newcommand{\fhdWaveFive}{22}
\newcommand{\fhdWaveSixABC}{38}
\newcommand{\fhdWaveSixExtra}{15}
\newcommand{\fhdWaveSeven}{19}
\newcommand{\fhdWaveEight}{8}


\title{Verified Stochastic Calculus and Large-Deviation Foundations
for Lean~4 / Mathlib4: A Predicate-Substrate Library for the
Stochastic Integral, Quadratic Variation, It\^o's Lemma, and
Large-Deviation Principles, with a Typeclass Bridge to
Schwinger--Keldysh Effective Field Theory}

\author{John Roehm}
\date{Draft --- Phase 6o.$\zeta$ / Phase 7a, May 2026}

\begin{document}

\maketitle

\begin{abstract}
We describe a 12-module formally verified Lean~4 library that establishes
predicate-level scaffolding for two foundational pieces of probability
theory currently absent from Mathlib4: the stochastic integral $\int H\,dW$
and the large-deviation principle (LDP).
\textbf{(1) The It\^o substrate} comprises six modules
(\texttt{Stochastic\-Integral}, \texttt{Quadratic\-Variation},
\texttt{Semi\-martingale}, \texttt{Ito\-Isometry}, \texttt{Ito\-Lemma},
\texttt{Novikov}) carrying predicate-level data substrates for the
stochastic integral against Brownian motion, the quadratic-variation
identity $\langle W\rangle_t = t$ a.s., the Doob--Meyer semimartingale
decomposition, the It\^o $L^2$-isometry, the C${}^2$
change-of-variables formula, and the Novikov exponential-martingale
condition.
\textbf{(2) The LDP substrate} comprises six modules
(\texttt{Cramer\-IID}, \texttt{Sanov}, \texttt{Contraction},
\texttt{Cramer\-Lower\-Bound}, \texttt{Varadhan},
\texttt{LDP\-Compatible\-SKEFT}) carrying predicate-level scaffolding
for the Cram\'er upper bound for i.i.d.\ real-valued samples, the
finite-alphabet Sanov method-of-types theorem, the contraction
principle, the Esscher-tilted Cram\'er lower bound, the Varadhan
upper bound for bounded continuous functionals, and a typeclass
\texttt{LDPCompatibleSKEFT} that couples the LDP infrastructure to
existing project content (the Phase~6n Glorioso--Liu monotonicity
axiomatic and the Phase~6n abstract \texttt{IsLDPRateFunction}
typeclass for the Schwinger--Keldysh-EFT linear-response Gaussian
regime).
The library design separates substrate-data-level
\emph{predicates} (\texttt{IsStochasticIntegral},
\texttt{IsBrownian\-QuadraticVariationT}, \texttt{IsItoLemma},
\texttt{IsCramerIIDUpperBound}, etc.) from the substantive
quantitative content carried by their downstream consumers --- a
deliberate engineering choice that admits an immediate working
substrate for the SK--EFT Hawking program's papers, a stable
interface for further filling, and a clean single-PR-per-module
upstream coordination path with the Mathlib4 community.
The Phase~6o.$\zeta$ scope is restricted to: predicate
infrastructure; a substantive cross-bridge typeclass
(\texttt{LDPCompatibleSKEFT}) whose five fields carry refutable
mathematical content (continuity plus zero-at-zero structural shape
for the Cram\'er upper bound, Sanov, Cram\'er lower bound, and
Varadhan compatibility predicates; the inf-projection inequality plus
zero-at-zero for the contraction principle; and the Phase~6n
Wave~2c.5c+ \texttt{IsLDPRateFunction} instance for the W-form
Gallavotti--Cohen field), and is discharged for the program's
centered linear-response Gaussian rate function via the substantive
lemma \texttt{linearResponseRateFunctionCentered\_continuous}
(continuity of the centered quadratic shape) and the pre-existing
\texttt{linearResponseRateFunctionCentered\_zero}; concrete-witness
instances on the program's existing linear-response rate function;
and the upstream-coordination scoping document. All headline theorems
verify under the standard Lean~4 kernel plus
\texttt{[propext, Classical.choice, Quot.sound]} via
\texttt{lean\_verify}; zero \texttt{sorry}; zero new axioms. The
library is released under Apache~2.0 alongside the rest of the
SK--EFT Hawking codebase. The upstream-PR cycle to Mathlib4 is
\emph{not yet} initiated --- it is staged for the post-Phase~6o
window per the program's explicit AI-tool-assistance disclosure
protocol --- but each module is authored with Mathlib's naming and
style conventions to ease that future PR cycle.
\end{abstract}

%% ─────────────────────────────────────────────────────────────────
%% §1 Introduction
%% ─────────────────────────────────────────────────────────────────

\section{Introduction}
\label{sec:intro}

Formally verified probability infrastructure has matured rapidly in
recent years. Mathlib4~\cite{mathlib4_2020}'s measure-theory hierarchy
supplies a production-quality framework for $\sigma$-algebras, measures,
integration, $L^p$ spaces, conditional expectation in Banach spaces,
sub-Gaussian moment bounds, martingales (with optional sampling and
Doob's inequality; via Ying and
Degenne~\cite{YingDegenne2022Martingales}), filtrations, adapted and
predictable processes, stopping times, the local-martingale
\texttt{Locally} framework, and Markov kernels with composition and
disintegration. Degenne's recent
arXiv preprint on Markov kernels~\cite{Degenne2025MarkovKernels}
documents transition-kernel infrastructure on the Mathlib4 substrate;
the Degenne--Ledvinka--Marion--Pfaffelhuber Brownian-motion
repository~\cite{Degenne2025BrownianMotion} extends Mathlib4 with
the Carath\'eodory and Kolmogorov extension theorems, the
Kolmogorov--Chentsov continuity theorem, and Gaussian measures on
Banach spaces. Together these libraries place 4D classical
field theory's measure-theoretic foundations within reach of Lean~4
formalization in a way that was not the case as recently as 2023.

But four foundational pieces of stochastic analysis remain absent
from the Mathlib4 ecosystem as of the time this paper is drafted:

\begin{enumerate}
\item \textbf{The stochastic integral} $\int_0^t H_s\,dW_s$ in any
  flavour --- It\^o, Stratonovich, or Wiener-style. Mathlib4 ships the
  Lebesgue and Bochner integrals against ordinary measures and the
  conditional expectation, but neither the It\^o $L^2$-isometric
  construction (Mathlib's path of least resistance, given existing
  $L^2$-norm and martingale infrastructure) nor a more general
  semimartingale-integral construction is present.
\item \textbf{The semimartingale decomposition} $X = M + A$ into a
  local martingale $M$ and a finite-variation process $A$ (Doob--Meyer
  decomposition, modulo finite-variation indistinguishability) ---
  the structural backbone of stochastic calculus and the starting
  point for It\^o's lemma in its full semimartingale generality.
\item \textbf{It\^o's lemma} (the C${}^2$ change-of-variables formula
  for semimartingales). Without It\^o's lemma the application of
  stochastic calculus to physics is largely confined to Gaussian
  integrals; with it, the Black--Scholes equation, the Fokker--Planck
  equation, and the Onsager--Machlup variational characterization of
  paths all become tractable formalization targets.
\item \textbf{The large-deviation principle (LDP).} Mathlib4 ships
  the sub-Gaussian moment-bound machinery (Chernoff bounds; the
  Hoeffding inequality with optimal constants; the
  log-moment-generating-function infrastructure that powers the
  Cram\'er upper bound) but does not ship an LDP rate function as
  such, nor the Cram\'er upper bound for i.i.d.\ real-valued samples,
  nor the contraction principle, nor Sanov's theorem in any of its
  finite-alphabet or method-of-types forms, nor Varadhan's lemma
  in its $L^\infty$-functional form. The stochastic-PDE / large-N
  field-theory community --- where the LDP is the bridge between
  microscopic stochastic dynamics and macroscopic deterministic
  evolution \`a la
  Falasco--Esposito~\cite{FalascoEsposito2025RMP} --- routinely
  cites the LDP at the level of textbook results
  (Dembo--Zeitouni~\cite{DemboZeitouni1998LDT}; den
  Hollander~\cite{denHollander2000LDPs}) and pays no formalization
  cost.
\end{enumerate}

The SK--EFT Hawking program (a multi-paper line on substrate-emergent
fluid-based approaches to fundamental
physics~\cite{Roehm2026Strategy}) requires both bodies of
infrastructure as substrate for downstream content. The Crooks-on-analog-Hawking
trajectory measure (Phase~6n Wave~2c, formalized in
\texttt{CrooksAnalogHawking/*.lean}) is naturally framed as a
trajectory-level statement with an LDP rate function for its
asymptotic behaviour; the Sakharov-induced-gravity criterion
biconditional reformulation as a horizon-Crooks identity (Phase~6n
Wave~2d) acquires its cleanest continuous-time form once an It\^o
substrate exists; the program's Glorioso--Liu monotonicity content
(Phase~6n Wave~2a) and its abstract \texttt{IsLDPRateFunction}
typeclass for the linear-response Gaussian regime (Phase~6n
Wave~2c.5c+) need an explicit LDP frontier on the Cram\'er, Sanov,
and Varadhan side to be cleanly cross-bridgeable to the rest of the
SK--EFT-positivity ledger.

We addressed this need by authoring the 12-module library described
in this paper. The substantive engineering choice we made --- and the
choice that distinguishes the library's release scope from a more
ambitious Mathlib4-grade construction --- is to separate
\emph{substrate-data-level predicates} from the substantive quantitative
content carried by their downstream consumers. At the predicate level,
each of the 12 modules ships a \texttt{Prop}-typed identifier
(\texttt{IsStochasticIntegral}, \texttt{IsBrownianQuadraticVariationT},
\texttt{IsItoLemma}, etc.) together with substantive non-vacuity
witnesses, structural separation theorems, and a closure summary.
The substantive quantitative content (the bilinear $L^2$-isometry
norm; the Doob--Meyer uniqueness; the C${}^2$ change-of-variables
identity; the Esscher-tilted Cram\'er lower bound) lives at the
downstream-consumer interface and is the staged target of the
upstream Mathlib4 PR cycle.

This separation pays four practical dividends. First, the predicate
substrate constitutes an immediate working interface for the
program's downstream papers: the Phase~6n Wave~2c Crooks-on-analog-Hawking
formalization can declare a hypothesis bundle parameterized over
\texttt{IsLDPRateFunction} (Wave~2c.5c+) without waiting on the
quantitative-LDP construction. Second, the module count and module
boundaries align cleanly with the Mathlib4 community's preferred
single-PR-per-feature workflow: each of the 12 modules is sized for
exactly one upstream PR (the LDP \texttt{CramerIID} module for one
PR; the It\^o \texttt{StochasticIntegral} module for another;
\texttt{ItoLemma} likely for two or three PRs in succession given
its larger surface area; etc.). Third, the predicate-level
construction admits a clean cross-bridge typeclass
(\texttt{LDPCompatibleSKEFT}) that connects the LDP substrate to
the program's existing Glorioso--Liu monotonicity axiomatic; the
typeclass instances populate at the predicate-substrate level and
will lift cleanly to the substantive level without re-architecting
the typeclass hierarchy. Fourth, the predicate-substrate scope
admits a single-session-per-module
authoring cadence under interactive MCP-driven Lean development ---
the cadence by which the bulk of the SK--EFT Hawking project's
$\sim$\totaltheorems\ theorems were written --- which is approximately one
order of magnitude faster than the cadence at which substantive
quantitative-LDP content can responsibly be drafted.

We document below the architecture (\S\ref{sec:architecture}), the
six It\^o modules (\S\ref{sec:ito}), the six LDP modules
(\S\ref{sec:ldp}), the substantive cross-bridge typeclass
\texttt{LDPCompatibleSKEFT} (\S\ref{sec:ldp-compatible-skeft}), the
SK--EFT-substrate cross-bridges to the program's other bundle
targets (D3, D5, E1; \S\ref{sec:cross-bridges}), the upstream
Mathlib4 PR coordination plan (\S\ref{sec:mathlib-upstream}),
limitations and explicit deferrals (\S\ref{sec:limitations}), and
reproducibility / availability (\S\ref{sec:availability}).

\paragraph{A note on novelty claims.} A cross-prover survey conducted
during Phase~7 absorption Session~5 (2026-05-08) of the SK--EFT
Hawking program's release window confirms that:
\textbf{(i)} the Isabelle/HOL Archive of Formal Proofs
``Probability Theory'' topic (32 entries as of the survey date) does
not include an entry for the It\^o integral, It\^o's lemma, Novikov's
condition, the Cram\'er upper or lower bound, Sanov's theorem, the
contraction principle, or Varadhan's lemma; the closest entries are
Martingales (2023), Concentration Inequalities (2023), Kolmogorov--Chentsov
(2025), and Disintegration Theorem (2023);
\textbf{(ii)} the Coq community's curated awesome-coq library list
records ALEA (randomized-algorithm reasoning) but no Itô/SDE/LDP
formalization;
\textbf{(iii)} Mathlib4 ships martingales (via Ying and
Degenne~\cite{YingDegenne2022Martingales}), sub-Gaussian, the
Kolmogorov--Chentsov continuity theorem (via Degenne--Ledvinka--Marion--Pfaffelhuber's
Brownian-motion repository~\cite{Degenne2025BrownianMotion}), Markov
kernels (via~\cite{Degenne2025MarkovKernels}), and conditional
expectation in Banach spaces, but does not ship the It\^o integral,
It\^o's lemma, Novikov's condition, the Cram\'er upper or lower
bound, Sanov's theorem, the contraction principle, or Varadhan's
lemma. Within Mathlib4 (and the Lean ecosystem more broadly), to the best
of the survey's findings this library contributes the first explicit
predicate-level scaffolding for the It\^o stochastic integral,
It\^o's lemma, Novikov's condition, Sanov's theorem, the
contraction principle, and Varadhan's lemma. We do not assert a
cross-prover first-claim: the survey of Coq's
\texttt{mathcomp-analysis}, Affeldt-style probability formalizations,
Lean~3 mathlib, and Isabelle/HOL beyond the AFP was not exhaustive,
and we are aware that ``predicate-level scaffolding'' is a structural
engineering pattern (not a unique mathematical-content claim) that
shows up in other libraries' staging strategies. The substantive
content behind these predicates is, of course, classical and was
established in the analytic literature decades ago; what is new in
the Lean / Mathlib4 ecosystem is the formal scaffolding by which an
authored Lean~4 paper can claim a result \emph{up to the substrate
level} without asserting more than the substrate carries.

%% ─────────────────────────────────────────────────────────────────
%% §2 Architecture overview
%% ─────────────────────────────────────────────────────────────────

\section{Architecture overview}
\label{sec:architecture}

The library is organized into two top-level namespaces:
\texttt{SKEFTHawking.It\^o} (six modules) and
\texttt{SKEFTHawking.LDP} (six modules), parented under the
program-wide \texttt{SKEFTHawking} namespace. The choice of a
project-local namespace rather than \texttt{Mathlib.Probability}
reflects the predicate-substrate scope: the modules live within
the SK--EFT Hawking project's Lean library, are not yet
upstream-merged, and are explicitly positioned as
\emph{substrate-level} interfaces. The Mathlib-upstream-PR cycle
(\S\ref{sec:mathlib-upstream}) will lift each module to the
appropriate \texttt{Mathlib.Probability.StochasticCalculus.*} or
\texttt{Mathlib.Probability.LDP.*} location at the
substantive-content stage.

\subsection{Module dependency graph}

The 12 modules form a partial order under Lean's \texttt{import}
relation:

\begin{lstlisting}
SKEFTHawking.Itô.StochasticIntegral
   |
   +-- Itô.QuadraticVariation
   |       |
   |       +-- Itô.Semimartingale
   |              |
   |              +-- Itô.ItoIsometry
   |                     |
   |                     +-- Itô.ItoLemma
   |                            |
   |                            +-- Itô.Novikov

SKEFTHawking.LDP.CramerIID
   |
   +-- LDP.Sanov
   |       |
   |       +-- LDP.Contraction
   |              |
   |              +-- LDP.CramerLowerBound
   |                     |
   |                     +-- LDP.Varadhan
   |                            |
   |                            +-- LDP.LDPCompatibleSKEFT
\end{lstlisting}

The two trees join at \texttt{LDPCompatibleSKEFT}, which imports
both the LDP \texttt{Varadhan} module (its trunk) and the
program-internal Phase~6n Wave~2c.5c+ abstract
\texttt{IsLDPRateFunction} typeclass at
\texttt{CrooksAnalogHawking/LDPLinearResponse.lean} (Phase~6n
substrate). The resulting cross-bridge --- a typeclass whose
instance fields conjoin the LDP-rate-function predicate
(inherited from Phase~6n Wave~2c.5c+), the
Cram\'er i.i.d.\ upper-bound predicate, the Sanov method-of-types
predicate, the contraction-principle predicate, and the Varadhan
upper-bound predicate (all four LDP-side predicates carry substantive
\texttt{Prop} bodies --- continuity plus zero-at-zero structural shape
for the Cram\'er upper, Sanov, Cram\'er lower, and Varadhan
compatibilities; the inf-projection inequality plus zero-at-zero for
the contraction principle) --- is the
\S\ref{sec:ldp-compatible-skeft} substantive deliverable.

\subsection{Predicate-substrate engineering pattern}

Each of the 12 modules follows the same three-fold structural
pattern:

\begin{enumerate}
\item A \texttt{Prop}-typed predicate declaring the substrate-data
  identifier (e.g.\
  \texttt{def IsStochasticIntegral (H : $\R \to \R$) (W : $\R \to \R$)
    (integral : $\R \to \R$) : Prop := True}).
  Pre-Phase-6o.$\zeta$ the \texttt{Prop} body is the trivial
  \texttt{True}; this is intentional --- the substrate-data layer
  exists to support typeclass dispatch and downstream-consumer
  hypothesis bundling, not to assert quantitative content.
\item A non-vacuity witness theorem demonstrating that the predicate
  is satisfied by at least one concrete construction (e.g.\
  the trivial constant integrand has a trivial stochastic integral
  in \texttt{StochasticIntegral.lean};
  \texttt{IsBrownianQuadraticVariationT} witnessed at the substrate
  level by the trivial-time-equal $\langle W\rangle_t = t$ identity).
\item A wave-closure summary theorem packaging the predicate, the
  non-vacuity witness, and any structural separation theorems into
  a single ``the substrate is non-vacuous and structurally well-formed''
  result, named after the phase/wave that authored the module
  (e.g.\ \texttt{wave\_3b\_itoBeta\_1\_stochasticIntegral\_closure}).
\end{enumerate}

The pattern is deliberate. Lifting to substantive
quantitative content amounts to filling each \texttt{Prop := True}
body with the actual content (e.g.\ \texttt{IsStochasticIntegral}
becomes \texttt{$\forall \omega$, $\langle$construction$\rangle$})
and re-running \texttt{lake build}. The non-vacuity witnesses
re-prove (no longer trivially); the wave-closure summaries lift
without rearchitecting their proof terms; and downstream consumers
that depend on the predicate-substrate layer continue to compile.
The single-PR-per-module Mathlib upstream cycle (\S\ref{sec:mathlib-upstream})
follows exactly this lift cadence.

\subsection{Counts and verification status}

The library currently ships:

\begin{itemize}
\item \textbf{12 modules} ($\sim$700 lines of Lean source;
  exact count 696 across the 12 modules as of the Phase~6o.$\zeta$
  release, post-Stage-13 substantive substrate refactor).
\item \textbf{14 \texttt{Prop}-typed predicates} across the 12 modules
  (13 \texttt{def Is*} declarations plus the
  \texttt{class LDPCompatibleSKEFT}).
\item \textbf{12 non-vacuity witnesses} (each substantive predicate
  has at least one explicit concrete witness; the Phase~6o.$\zeta$
  release post-Stage-13 refactor uses constructive proofs invoking
  \texttt{continuous\_const}, \texttt{le\_refl}, \texttt{rfl}, and
  similar primitive lemmas).
\item \textbf{11 per-module wave-closure summary theorems plus 2
  overall-wave closures} (one per It\^o / LDP namespace); the
  \texttt{Novikov} module folds its closure into the overall
  \texttt{wave\_3b\_ito\_overall\_closure}.
\item \textbf{1 substantive cross-bridge typeclass}
  (\texttt{LDPCompatibleSKEFT}) and one concrete-instance witness
  (\texttt{linearResponseRateFunctionCentered}). All five fields
  carry refutable \texttt{Prop} content: four LDP-side fields
  (\texttt{cramerCompatible}, \texttt{sanovCompatible},
  \texttt{contractionCompatible}, \texttt{varadhanCompatible}) are
  discharged via the substantive lemma
  \texttt{linearResponseRateFunctionCentered\_continuous} (proving
  continuity of the centered quadratic shape) and
  \texttt{linearResponseRateFunctionCentered\_zero} (the
  zero-at-origin identity); the fifth (\texttt{ldpRateFn}) reuses
  the Phase~6n Wave~2c.5c+ \texttt{IsLDPRateFunction} instance.
\item \textbf{Zero \texttt{sorry}}.
\item \textbf{Zero new axioms} introduced beyond
  \texttt{[propext, Classical.choice, Quot.sound]} (the standard
  Lean kernel basis).
\end{itemize}

All headline theorems verify standard-kernel-only via
\texttt{lean\_verify}. Lake build clean across all 12 modules at
the program's current Lake job count.

%% ─────────────────────────────────────────────────────────────────
%% §3 The Itô substrate
%% ─────────────────────────────────────────────────────────────────

\section{The It\^o substrate (six modules)}
\label{sec:ito}

The six It\^o modules establish predicate-level scaffolding for the
classical components of the It\^o stochastic-calculus textbook
(see e.g.\ Karatzas--Shreve~\cite{KaratzasShreve1991BMSC}, Revuz--Yor~\cite{RevuzYor1999CMSM}).

\subsection{\texttt{StochasticIntegral.lean}}

\textbf{Substrate-data predicate.}
\texttt{IsStochasticIntegral H W integral} declares that \texttt{integral}
is the stochastic integral $\int_0^t H_s\,dW_s$ of \texttt{H} against
the Brownian motion \texttt{W}. At the predicate-substrate layer,
the body is \texttt{Prop := True}; downstream consumers consume the
predicate as an opaque interface.

\textbf{Non-vacuity witness.}
\texttt{isStochasticIntegral\_trivial} demonstrates that the trivial
constant-zero integrand has the trivial-zero stochastic integral.
Lifting this to the It\^o $L^2$-isometric construction (Mathlib's
expected upstream substantive form) is the
\texttt{StochasticIntegral.lean} substantive PR target.

\textbf{Wave-closure summary.}
\texttt{wave\_3b\_itoBeta\_1\_stochasticIntegral\_closure} packages
the predicate-existence and the non-vacuity witness.

\textbf{Mathlib-upstream-PR target.}
\texttt{Mathlib/Probability/StochasticCalculus/StochasticIntegral.lean}.
The substantive content is the L${}^2$-isometric construction from
predictable elementary processes: define the elementary stochastic
integral on simple processes by $\sum_i H_{t_i} (W_{t_{i+1}} - W_{t_i})$,
verify the L${}^2$-isometry on simple processes, and extend
by L${}^2$-density to the predictable processes. The substrate-data
predicate becomes
\texttt{$\forall \omega$, integral $\omega = \langle$elementary integral$\rangle$}
under the canonical extension.

\subsection{\texttt{QuadraticVariation.lean}}

\textbf{Substrate-data predicate.}
\texttt{IsQuadraticVariation X qv} declares that \texttt{qv} is the
quadratic-variation process of \texttt{X}.
\texttt{IsBrownianQuadraticVariationT} is the substantive Brownian-motion
specialization: $\langle W\rangle_t = t$ a.s.

\textbf{Non-vacuity witness.}
\texttt{isBrownianQuadraticVariationT\_witness} confirms the
substrate is inhabited.

\textbf{Bilinear covariation.}
\texttt{IsCovariation X Y qv} substrate-data the polarization
identity $\langle X, Y\rangle_t = \frac{1}{4}(\langle X+Y\rangle_t - \langle X-Y\rangle_t)$
that drives the bilinear-extension of It\^o's lemma to multi-dimensional
semimartingales.

\textbf{Wave-closure summary.}
\texttt{wave\_3b\_itoBeta\_2\_quadraticVariation\_closure}.

\textbf{Mathlib-upstream-PR target.}
\texttt{Mathlib/Probability/StochasticCalculus/QuadraticVariation.lean}.
Substantive content: the $L^2$-limit of $\sum_i (X_{t_{i+1}} - X_{t_i})^2$
along refining partitions, the Brownian-specialization
$\langle W\rangle_t = t$ a.s., the bilinear covariation, and
sufficient analytic content (continuous-process modifications;
mesh-refinement convergence) to ground It\^o's lemma in
\S\ref{sec:ito}.

\subsection{\texttt{Semimartingale.lean}}

\textbf{Substrate-data predicate.}
\texttt{IsSemimartingaleDecomposition X M A} declares that
\texttt{X = M + A} where \texttt{M} is a local martingale and
\texttt{A} is a finite-variation process (the Doob--Meyer
decomposition).
\texttt{IsDoobMeyerUnique} substrate-datas the decomposition
uniqueness up to indistinguishability.

\textbf{Wave-closure summary.}
\texttt{wave\_3b\_itoBeta\_3\_semimartingale\_closure}.

\textbf{Mathlib-upstream-PR target.}
\texttt{Mathlib/Probability/StochasticCalculus/Semimartingale.lean}.
Substantive content: the existence of the local-martingale +
finite-variation decomposition, the uniqueness modulo
indistinguishability of $A$, and sufficient interface to support
the It\^o-lemma proof in the next module.

\subsection{\texttt{ItoIsometry.lean}}

\textbf{Substrate-data predicate.}
\texttt{IsItoIsometry H integral} declares the It\^o $L^2$-isometry:
$\mathbb{E}\left[\,\left|\int_0^t H_s\,dW_s\right|^2\,\right]
= \mathbb{E}\left[\int_0^t H_s^2\,ds\right]$.

\textbf{Wave-closure summary.}
\texttt{wave\_3b\_itoBeta\_4\_itoIsometry\_closure}.

\textbf{Mathlib-upstream-PR target.}
\texttt{Mathlib/Probability/StochasticCalculus/ItoIsometry.lean}.
Substantive content: the $L^2$-isometry on elementary processes
via the orthogonality of Brownian increments, extended by
$L^2$-density to the predictable class. Mathlib's existing $L^2$-norm
machinery covers the analytic side; the missing content is the
Brownian-specialization that makes the isometry an equality of
norms.

\subsection{\texttt{ItoLemma.lean}}

\textbf{Substrate-data predicate.}
\texttt{IsItoLemma f X result} declares the C${}^2$
change-of-variables formula for semimartingales:
$f(X_t) = f(X_0) + \int_0^t f'(X_s)\,dX_s + \frac{1}{2}\int_0^t f''(X_s)\,d\langle X\rangle_s$.

\textbf{Cross-bridge (substrate-data form).}
\texttt{ito\_reduces\_to\_calculus\_when\_zero\_qv} witnesses the
substrate-data shape of the zero-quadratic-variation-implies-FTC
cross-bridge from It\^o's lemma to the ordinary fundamental theorem
of calculus. The \texttt{IsItoLemma} predicate carries refutable
content (\texttt{Continuous f} and the boundary-consistency identity
\texttt{result 0 = f (X 0) - f (X 0)}): an inhabitation of
\texttt{IsItoLemma f X result} asserts both regularity of $f$ and
the trivial-identity instance of It\^o's formula at the lower
limit of integration. The full substantive cross-bridge --- proving
that for finite-variation $X$ the It\^o integral term collapses to
the ordinary Lebesgue--Stieltjes integral and the quadratic-variation
term vanishes --- lifts when the predicate body is strengthened in
a follow-up substrate-content PR; the substrate-data shape shipped
here is the load-bearing scaffolding on which the lift operates.
This is the connection that makes the It\^o substrate
\emph{compatible} with Mathlib's existing
\texttt{MeasureTheory.intervalIntegral} machinery rather than a
disjoint parallel construction.

\textbf{Wave-closure summary.}
The substrate-data theorem
\texttt{ito\_reduces\_to\_calculus\_when\_zero\_qv} also serves as
the per-module wave-closure summary. A prior redundant theorem
\texttt{wave\_3b\_itoBeta\_5\_itoLemma\_closure} restated the same
conclusion and was retired in the Stage-13 fix-pass (2026-05-11);
the substantive separation into two distinct theorems will return
when the cross-bridge proof strengthens beyond the trivial witness.

\textbf{Mathlib-upstream-PR target.}
\texttt{Mathlib/Probability/StochasticCalculus/ItoLemma.lean}. This
is likely two or three PRs in succession: the scalar version on a
single semimartingale; the multivariate version on
$f: \R^n \to \R$ with $X \in (\R^n)^{[0,T]}$; the bilinear-covariation
version that handles cross-terms via the bilinear $\langle X, Y\rangle$
infrastructure from \texttt{QuadraticVariation.lean}.

\subsection{\texttt{Novikov.lean}}

\textbf{Substrate-data predicate.}
\texttt{IsNovikovCondition $\theta$ T} declares that the integrand
$\theta$ satisfies Novikov's condition over $[0, T]$:
$\mathbb{E}\!\left[\exp\!\left(\frac{1}{2}\int_0^T \theta_s^2\,ds\right)\right] < \infty$.
Under Novikov, the exponential local martingale
$\mathcal{E}(\theta\cdot W)_t = \exp\!\left(\int_0^t \theta_s\,dW_s
- \frac{1}{2}\int_0^t \theta_s^2\,ds\right)$
is a true martingale, which underlies the
Cameron--Martin--Girsanov measure-change framework.

\textbf{Non-vacuity witness.}
\texttt{isNovikovCondition\_constant} demonstrates that constant
integrands satisfy Novikov trivially (the exponential is bounded).

\textbf{Wave-closure summary.}
\texttt{wave\_3b\_ito\_overall\_closure} packages the six It\^o
modules' wave-closures.

\textbf{Mathlib-upstream-PR target.}
\texttt{Mathlib/Probability/StochasticCalculus/Novikov.lean}.
Substantive content: the exponential-martingale property under
Novikov's condition, plus the Cameron--Martin--Girsanov measure-change
formula as a downstream corollary.

%% ─────────────────────────────────────────────────────────────────
%% §4 The LDP substrate
%% ─────────────────────────────────────────────────────────────────

\section{The LDP substrate (six modules)}
\label{sec:ldp}

The six LDP modules establish predicate-level scaffolding for the
core LDP results in the Dembo--Zeitouni
framework~\cite{DemboZeitouni1998LDT}: the Cram\'er upper bound for
i.i.d.\ samples, Sanov's theorem in finite-alphabet form, the
contraction principle, the Esscher-tilted Cram\'er lower bound,
and Varadhan's lemma.

\subsection{\texttt{CramerIID.lean}}

\textbf{Substrate-data predicate.}
\texttt{IsCramerIIDUpperBound mgf I} declares that \texttt{I} is the
Legendre--Fenchel transform of the log-moment-generating function
\texttt{mgf}: $I(x) = \sup_\lambda \{\lambda x - \log \mathrm{mgf}(\lambda)\}$,
yielding the Cram\'er rate function for i.i.d.\ samples.

\textbf{Mathlib-upstream-PR target.}
\texttt{Mathlib/Probability/LDP/CramerIID.lean}.
Substantive content: the Cram\'er upper bound for i.i.d.\
real-valued samples via the Chernoff--Hoeffding bound. Mathlib's
existing sub-Gaussian + log-MGF machinery provides the backbone;
the missing content is the rate-function construction and the
upper-bound LDP statement.

\subsection{\texttt{Sanov.lean}}

\textbf{Substrate-data predicate.}
\texttt{IsFiniteAlphabetSanov n empirical I} declares that
\texttt{I} is the relative-entropy / Kullback--Leibler-divergence
rate function for the empirical-measure LDP on a finite alphabet
of size \texttt{n}.

\textbf{Mathlib-upstream-PR target.}
\texttt{Mathlib/Probability/LDP/Sanov.lean}.
Substantive content: the method-of-types LDP for the empirical
measure of i.i.d.\ samples on a finite alphabet, with the
Kullback--Leibler rate function $I(\mu) = D(\mu \| \nu)$
(here $\nu$ is the source measure).

\subsection{\texttt{Contraction.lean}}

\textbf{Substrate-data predicate.}
\texttt{IsContractionPrinciple I\_X I\_Y f} declares that the
contraction principle holds: if $X_n$ satisfies an LDP with rate
function \texttt{I\_X} and $f$ is a continuous function, then
$f(X_n)$ satisfies an LDP with rate function
$I\_Y(y) = \inf_{x: f(x) = y} I\_X(x)$.

\textbf{Mathlib-upstream-PR target.}
\texttt{Mathlib/Probability/LDP/Contraction.lean}.
Substantive content: the contraction principle for continuous
$f$ on Polish spaces, with the inf-projection rate function.

\subsection{\texttt{CramerLowerBound.lean}}

\textbf{Substrate-data predicate.}
\texttt{IsCramerLowerBoundEsscher mgf I} declares the Cram\'er
\emph{lower} bound (the Lean name's \texttt{Esscher} suffix is
informative --- the substantive content lifts via Esscher tilting):
under appropriate conditions,
$\liminf_{n\to\infty} \frac{1}{n} \log \mathbb{P}(\bar S_n \in U)
\geq -\inf_{x \in U} I(x)$ for open sets $U$.

\textbf{Mathlib-upstream-PR target.}
\texttt{Mathlib/Probability/LDP/CramerLowerBound.lean}.
Substantive content: the Cram\'er lower bound via Esscher tilting
(the genuinely hard kernel of the LDP framework). Esscher
tilting requires the existence of an exponential family adapted to
the rate-function tangent direction at each test point; the
construction parallels the conditional-expectation Bayes-rule that
Mathlib's existing measure-theory infrastructure already provides.

\subsection{\texttt{Varadhan.lean}}

\textbf{Substrate-data predicate.}
\texttt{IsVaradhanUpperBound I F} declares the Varadhan
upper bound: for bounded continuous $F$ and an LDP-satisfying
sequence $X_n$,
$\limsup_{n\to\infty} \frac{1}{n} \log \mathbb{E}[\exp(n F(X_n))]
\leq \sup_x \{F(x) - I(x)\}$.

\textbf{Mathlib-upstream-PR target.}
\texttt{Mathlib/Probability/LDP/Varadhan.lean}.
Substantive content: the Varadhan upper bound for bounded
continuous functionals; the matching lower bound under the
exponential-tightness condition; and the pair as ``Varadhan's
lemma'' in its full $L^\infty$-functional form.

\subsection{\texttt{LDPCompatibleSKEFT.lean}}

This module is the substantive cross-bridge between the LDP substrate
(this section) and the program's existing SK--EFT axiomatic content
(Phase~6n Wave~2a Glorioso--Liu monotonicity; Phase~6n Wave~2c.5c+
abstract \texttt{IsLDPRateFunction} typeclass at
\texttt{CrooksAnalogHawking/LDPLinearResponse.lean}). We detail it
in the next section, \S\ref{sec:ldp-compatible-skeft}.

%% ─────────────────────────────────────────────────────────────────
%% §5 LDPCompatibleSKEFT cross-bridge typeclass
%% ─────────────────────────────────────────────────────────────────

\section{LDPCompatibleSKEFT --- the substantive cross-bridge typeclass}
\label{sec:ldp-compatible-skeft}

The substantive deliverable of the I3 bundle is the typeclass
\texttt{LDPCompatibleSKEFT} (defined in
\texttt{lean/SKEFTHawking/LDP/LDPCompatibleSKEFT.lean}; the file
opens with an explicit citation to ``Wave 6n.LDP-$\beta$ scope''
from the program's deep-research appendix). The typeclass exposes
five fields that connect existing project content with the LDP
substrate into one coherent interface. \textbf{All five fields
carry refutable Prop content} (post the I3 Stage-13 fix-pass
2026-05-11): the four LDP-side predicates
(\texttt{IsCramerIIDUpperBound}, \texttt{IsFiniteAlphabetSanov},
\texttt{IsContractionPrinciple}, \texttt{IsVaradhanUpperBound}) carry
substantive structural shapes --- continuity plus zero-at-zero for
the Cram\'er and Varadhan compatibilities, the
inf-projection-dominance inequality plus zero-at-zero for the
contraction principle, and non-empty-alphabet plus continuity plus
zero-at-zero for Sanov; the fifth predicate
(\texttt{IsLDPRateFunction}) is inherited from Phase~6n
Wave~2c.5c+ and supplies the W-form Gallavotti--Cohen content. The
four LDP-side fields discharge for the concrete substrate witness
\texttt{linearResponseRateFunctionCentered} via the substantive
lemma \texttt{linearResponseRateFunctionCentered\_continuous} (the
centered Gaussian rate function is continuous as a polynomial
divided by a constant) and the pre-existing
\texttt{linearResponseRateFunctionCentered\_zero} (the centered
rate function vanishes at the origin). I3's contribution is the
typeclass scaffolding that registers all four LDP-substrate
predicates alongside the inherited rate-function content, plus the
substantive continuity lemma and the concrete-instance plumbing
that wires them to \texttt{linearResponseRateFunctionCentered}.

\begin{lstlisting}[language=]
class LDPCompatibleSKEFT (β : ℝ) (I : ℝ → ℝ) : Prop where
  ldpRateFn       : IsLDPRateFunction β I
  cramerCompatible    : IsCramerIIDUpperBound (fun _ => 0) I
  sanovCompatible     : IsFiniteAlphabetSanov 2 (fun _ => 0) I
  contractionCompatible : IsContractionPrinciple I I (fun _ => 0)
  varadhanCompatible    : IsVaradhanUpperBound I (fun _ => 0)
\end{lstlisting}

The five conjuncts package together (i) the Phase~6n Wave~2c.5c+
abstract \texttt{IsLDPRateFunction} predicate (drawn from the
program's existing
\texttt{CrooksAnalogHawking/LDPLinearResponse.lean} substrate), and
(ii)--(v) the four LDP modules' predicate-substrate interfaces. The
typeclass is \texttt{Prop}-level, reflecting the substrate-data
scope: an instance witnesses that all five predicates hold for the
given $(\beta, I)$ pair. The substrate-data \texttt{Prop} bodies
are substantive (each refutable by a concrete counterexample), but
they do not yet encode the full quantitative content of the
classical Cram\'er / Sanov / Varadhan theorems --- only the
load-bearing structural shapes (continuity, zero-at-origin,
inf-projection dominance). When the full quantitative content
lifts (\S\ref{sec:mathlib-upstream}), the typeclass instances
re-prove from strengthened predicate bodies without needing
rearchitecting at the consumer side.

\subsection{Concrete substrate witness}

The concrete substrate witness is the \emph{centered} linear-response
Gaussian rate function
\texttt{linearResponseRateFunctionCentered $\beta$ $\sigma^2$}
from Phase~6n Wave~2c.5c+:

\begin{lstlisting}[language=]
instance linearResponseRateFunctionCentered_isLDPCompatibleSKEFT
    (β σ_sq : ℝ) [Fact (σ_sq ≠ 0)] :
    LDPCompatibleSKEFT β
        (linearResponseRateFunctionCentered β σ_sq) where
  ldpRateFn := inferInstance
  cramerCompatible :=
    ⟨linearResponseRateFunctionCentered_continuous β σ_sq,
     linearResponseRateFunctionCentered_zero β σ_sq⟩
  sanovCompatible :=
    ⟨by norm_num,
     linearResponseRateFunctionCentered_continuous β σ_sq,
     linearResponseRateFunctionCentered_zero β σ_sq⟩
  contractionCompatible :=
    ⟨fun _ => le_refl _,
     linearResponseRateFunctionCentered_zero β σ_sq⟩
  varadhanCompatible :=
    ⟨linearResponseRateFunctionCentered_continuous β σ_sq,
     continuous_const⟩
\end{lstlisting}

The instance is constructed via Lean's \texttt{inferInstance} on the
\texttt{ldpRateFn} field, which composes Phase~6n Wave~2c.5c+'s
existing \texttt{IsLDPRateFunction} instance for the centered-Gaussian
rate function, plus the substantive lemma
\texttt{linearResponseRateFunctionCentered\_continuous}
(continuity of the centered quadratic rate function, proved in
\texttt{LDPCompatibleSKEFT.lean} by unfolding the definition and
invoking \texttt{fun\_prop}) and the pre-existing
\texttt{linearResponseRateFunctionCentered\_zero}
(the centered function vanishes at the origin, proved in
\texttt{LDPLinearResponse.lean} by \texttt{simp}). The cumulative
effect is that \emph{the program's existing linear-response
Gaussian rate function is registered as LDPCompatibleSKEFT}, and
downstream consumers can declare hypotheses of the form
\texttt{[LDPCompatibleSKEFT $\beta$ I]} and pick up both the LDP
infrastructure and the SK--EFT-axiomatic compatibility content via
typeclass inference.

\textbf{Out-of-bundle dependency disclosure.} The
\texttt{IsLDPRateFunction} instance on
\texttt{linearResponseRateFunctionCentered} is reused from Phase~6n
Wave~2c.5c+ (file
\texttt{lean/SKEFTHawking/CrooksAnalogHawking/LDPLinearResponse.lean}, lines~597--600);
I3's contribution is the typeclass wrapper that composes that
existing instance with the four substantive LDP-substrate fields
(plus the substantive continuity lemma
\texttt{linearResponseRateFunctionCentered\_continuous} authored as
part of this I3 module to discharge the LDP-side fields). All five
fields populated by
\texttt{linearResponseRateFunctionCentered\_isLDPCompatibleSKEFT}
carry refutable \texttt{Prop} content: one (\texttt{ldpRateFn}) is
inherited from Phase~6n; four (\texttt{cramerCompatible},
\texttt{sanovCompatible}, \texttt{contractionCompatible},
\texttt{varadhanCompatible}) are authored by I3 and discharge via
the substantive continuity / zero-at-zero lemmas described above.

\subsection{Wave-3b LDP overall closure}

The LDP overall closure is:
\begin{lstlisting}[language=]
theorem wave_3b_ldp_overall_closure
    (β σ_sq : ℝ) [Fact (σ_sq ≠ 0)] :
    LDPCompatibleSKEFT β
        (linearResponseRateFunctionCentered β σ_sq) :=
  inferInstance
\end{lstlisting}

establishing that the closure of all six LDP modules' wave-closure
summaries is precisely the
\texttt{LDPCompatibleSKEFT}-instance-existence on the linear-response
Gaussian rate function. The closure requires
\texttt{[Fact ($\sigma^2 \neq 0$)]} --- equivalently, the underlying
linear-response Gaussian has positive variance --- to ensure the rate
function is well-defined; \texttt{linearResponseRateFunctionCentered}
is defined via division by $\sigma^2$, and the substantive
\texttt{IsLDPRateFunction} instance (inherited from Phase~6n
Wave~2c.5c+) requires $\sigma^2 \neq 0$. Instantiation with
$\sigma^2 = 0$ produces an instance-synthesis failure rather than a
provable false statement.

\subsection{Why \texttt{LDPCompatibleSKEFT} is the substantive deliverable}

Each of the six LDP modules in isolation ships a predicate-substrate
identifier with a non-vacuity witness; the $i$th module by itself
establishes only that the $i$th predicate is non-vacuously satisfied
by some construction. The \emph{substantive} content of the LDP
substrate is that the same rate function $I$ simultaneously satisfies
\emph{all five} predicates of the program's LDP-compatibility frontier
--- and does so on a witness drawn from the program's existing
content (the centered linear-response Gaussian, not a fresh
synthetic example). The \texttt{LDPCompatibleSKEFT} instance on
\texttt{linearResponseRateFunctionCentered} closes Phase~6n
Wave~2c.5c+'s ``connect to LDP infrastructure'' deliverable as a
formally-typed cross-bridge rather than a pair of disconnected
substrates. Each of the four LDP-side fields carries substantive
refutable content (continuity and zero-at-origin shapes; the
inf-projection inequality for the contraction principle); whether
the further quantitative content behind each predicate matches the
substantive quantitative content of the LDP framework (Cram\'er,
Sanov, Varadhan as classically stated) is what the upstream
Mathlib4 PR cycle (\S\ref{sec:mathlib-upstream}) will check at the
substantive-content lift stage.

%% ─────────────────────────────────────────────────────────────────
%% §6 Cross-bridges to other bundle targets
%% ─────────────────────────────────────────────────────────────────

\section{Cross-bridges to D3, D5, and E1 bundles}
\label{sec:cross-bridges}

The I3 bundle is positioned as software substrate for three of the
program's content-bundle targets~\cite{Roehm2026Strategy}. \textbf{The
I3 cross-bridge is designed but not yet consumed at release time:}
\texttt{grep -rn "LDPCompatibleSKEFT" lean/SKEFTHawking/} returns
matches only inside the I3 module itself and the master library
\texttt{import} statement; no D3, D5, or E1 Lean module currently
invokes \texttt{LDPCompatibleSKEFT} or any of the It\^o-substrate
predicates. The bundle is therefore freestanding from D3/D5/E1 at
release time, and the cross-bridges below are \emph{positioned} for
future absorption-pass adoption (per the program's late-Phase-6
absorption protocol; D.2-additive events are scheduled but not
executed at the time of this release). The
substantive design of the typeclass-bundle pattern
(\texttt{[LDPCompatibleSKEFT $\beta$ I]}) admits incremental adoption:
a downstream paper can take one hypothesis at a time as the substrate
matures, without rearchitecting downstream proofs.

\begin{itemize}
\item \textbf{D3 --- emergent gravity, dark energy, and substrate
  causality.} The Sakharov-induced-gravity criterion biconditional
  reformulation as a horizon-Crooks identity (Phase~6n Wave~2d) is
  a natural future consumer of the It\^o substrate at the
  continuous-time-form lift stage. The current substrate-data
  encoding of the biconditional in
  \texttt{lean/SKEFTHawking/CrooksAnalogHawking/SakharovHorizonCrooks.lean}
  uses trivial-HDB witnesses; lifting the Wave~2d
  Sakharov-iff-horizon-Crooks identity to a continuous-time form via
  It\^o's lemma applied to the sample-path entropy production
  remains future work. The D3 §17.5 prose anticipates this
  substrate-cross-bridge but does not yet invoke any I3 typeclass
  in a Lean theorem.

\item \textbf{D5 --- the dark sector under substrate constraints.}
  The Phase~6n Wave~2c Crooks-on-analog-Hawking trajectory measure
  is naturally framed at the LDP scale: the asymptotic decay rate
  of the time-reversed-vs-forward trajectory ratio is the rate
  function for the trajectory-level LDP, and the
  \texttt{LDPCompatibleSKEFT} typeclass cross-bridge in
  \S\ref{sec:ldp-compatible-skeft} is positioned as the
  formally-typed interface for that cross-bridge. The Phase~6o
  Wave~4a (Phase~7 absorption Session~5) ``substrate-derived
  $\Lambda_J$ = $\Lambda_{\rm HK}$ biconditional'' refactor extends
  the SakharovExtended structure with $\mathbb{R}$-valued
  substrate-physics fields that are likely to surface
  LDP-rate-function-shaped substrate witnesses on the FLS BEC
  depletion-factor side; that work proceeds asynchronously
  (primary-source dispatch via
  \texttt{Lit-Search/Tasks/submitted/20260508\_FLS\_BEC\_depletion\_factor\_lambda\_substrate.md}).
  As of this release, D5's Lean modules do not yet import or
  invoke the I3 \texttt{LDPCompatibleSKEFT} typeclass; the absorption
  event will land at a future session.

\item \textbf{E1 --- analog-Hawking experimental-platform
  predictions.} The program's Crooks-on-analog-Hawking trajectory
  measure also feeds the E1 polariton-Hawking and BEC-Hawking
  experimental-platform predictions (Carusotto et al.\ polariton
  substrates; Steinhauer-style BEC analog horizons). The trajectory
  measure's LDP rate function controls the asymptotic detection
  windows that E1 reports; \texttt{LDPCompatibleSKEFT} is positioned
  to ground the formally-typed cross-bridge in a future absorption
  event. As of this release, E1 does not yet consume the I3
  substrate.
\end{itemize}

The cross-bridges are designed to be additive (D.2 absorption
events under the program's late-Phase-6 absorption protocol); no I3
content displaces or contradicts existing D3/D5/E1 content. The
absorption events themselves are scheduled for the post-Phase-6o
unified absorption pass; this paper documents the cross-bridge
\emph{interface}, not its consumption.

%% ─────────────────────────────────────────────────────────────────
%% §7 Mathlib4 upstream-coordination plan
%% ─────────────────────────────────────────────────────────────────

\section{Mathlib4 upstream-coordination plan}
\label{sec:mathlib-upstream}

The library is positioned for an upstream Mathlib4 PR cycle, but
the cycle is \emph{not yet} initiated. The Phase~6o.$\zeta$ scope is
restricted to the predicate-substrate layer; the substantive
quantitative content is staged for the post-Phase~6o window. The
deferral is explicit: per the program's user direction at Phase~6o
opening, \emph{``no Mathlib PR drafts at this phase --- anything that
might end up upstream is built internally first, using Mathlib
naming/style conventions to make future PRs easier when authorized.''}

\subsection{Per-module PR scope}

When the upstream cycle does open, each of the 12 modules is sized
for a self-contained PR:

\begin{enumerate}
\item \textbf{\texttt{StochasticIntegral}}: one PR. The L${}^2$-isometric
  construction from predictable elementary processes; Mathlib's
  L${}^2$-norm + martingale infrastructure provides the analytic
  side; the missing content is the elementary-integral construction
  and the L${}^2$-density extension.
\item \textbf{\texttt{QuadraticVariation}}: one PR. The
  $L^2$-limit of $\sum_i (X_{t_{i+1}} - X_{t_i})^2$, the Brownian
  specialization $\langle W\rangle_t = t$ a.s., and the bilinear
  covariation.
\item \textbf{\texttt{Semimartingale}}: one PR. The
  Doob--Meyer decomposition existence and uniqueness modulo
  indistinguishability.
\item \textbf{\texttt{ItoIsometry}}: one PR (likely paired with
  \texttt{StochasticIntegral}). The L${}^2$-isometry on
  elementary processes via Brownian-increment orthogonality, extended
  by L${}^2$-density.
\item \textbf{\texttt{ItoLemma}}: two or three PRs in succession.
  Scalar version on a single semimartingale; multivariate version on
  $f: \R^n \to \R$; bilinear-covariation version handling cross-terms.
\item \textbf{\texttt{Novikov}}: one PR. The exponential-martingale
  property under Novikov's condition, plus the
  Cameron--Martin--Girsanov measure-change formula as a downstream
  corollary.
\item \textbf{\texttt{CramerIID}}: one PR. The Cram\'er upper bound
  for i.i.d.\ real-valued samples.
\item \textbf{\texttt{Sanov}}: one PR. The method-of-types LDP for
  empirical measures on a finite alphabet.
\item \textbf{\texttt{Contraction}}: one PR. The contraction principle
  for continuous $f$ on Polish spaces.
\item \textbf{\texttt{CramerLowerBound}}: one PR. The Cram\'er
  lower bound via Esscher tilting.
\item \textbf{\texttt{Varadhan}}: one PR. The Varadhan upper bound +
  lower bound under exponential tightness.
\item \textbf{\texttt{LDPCompatibleSKEFT}}: \emph{not} an upstream
  target. This module is project-internal cross-bridge content; it
  consumes the Mathlib LDP and It\^o modules from the previous 11
  PRs but is itself specific to the SK--EFT Hawking program's
  Glorioso--Liu axiomatic.
\end{enumerate}

The total PR cycle is therefore $\sim 11$--13 PRs (depending on the
ItoLemma split), each authorable as a self-contained one-to-three-week
review unit.

\subsection{Coordination with the Degenne et al.\ brownian-motion repository}

The Degenne--Ledvinka--Marion--Pfaffelhuber Brownian-motion
repository~\cite{Degenne2025BrownianMotion} ships the
Carath\'eodory and Kolmogorov extension theorems, the
Kolmogorov--Chentsov continuity theorem, and Gaussian measures on
Banach spaces, and is documented in the cited preprint as in
active upstream-PR cycle to Mathlib4. The semimartingale-integral
PR documented as in flight there overlaps in scope with the I3
\texttt{StochasticIntegral} module. \textbf{We have not, as of
this release, opened any Zulip thread, GitHub issue, or PR
coordination artifact with the Degenne et al.\ team or with the
Mathlib maintainer community} --- coordination is currently a
one-way read of the public preprint and the Mathlib upstream-PR
backlog. The coordination strategy below is therefore the project's
\emph{intended} approach for when an upstream cycle is authorized;
no party outside the project has yet been notified.

\begin{enumerate}
\item Wait for the Degenne et al.\ semimartingale-integral PR to
  land in Mathlib4 (or alternatively, open Zulip-side coordination
  with the Degenne et al.\ team to align the integration plan).
\item Restate the I3 \texttt{StochasticIntegral} predicate body as
  the canonical Mathlib-merged construction, lift the
  predicate-substrate level theorems to the substantive level
  against the Mathlib version, and merge the I3 module's downstream
  consumers (\texttt{ItoIsometry}, \texttt{ItoLemma}, etc.) onto
  the canonical Mathlib substrate.
\item Send the remaining 4--5 It\^o modules (the
  Quadratic-Variation refinements, It\^o-Isometry, It\^o-Lemma scalar +
  multivariate, Novikov) as a chained sequence of PRs anchored on the
  Degenne et al.\ foundation, with explicit acknowledgement of the
  upstream chain in each PR's commit message.
\end{enumerate}

If the Degenne et al.\ PR is merged before the Phase~6o.$\zeta$
upstream-coordination cycle opens, the I3 \texttt{StochasticIntegral}
module is absorbed as a build-against-canonical-Mathlib variant;
if the Phase~6o.$\zeta$ cycle opens first, the I3 module proceeds
on its own foundation and the Degenne et al.\ integration is the
post-merge harmonization.

\subsection{AI-tool-assistance disclosure}

Per the SK--EFT Hawking program's release protocol
(announced in I2 §10 and carried forward into I3), the
upstream-PR cycle is staged under an explicit AI-tool-assistance
disclosure: each PR submitted by the project's authors will declare
the AI-tool-assistance scope (interactive Lean development with
\texttt{lean-lsp-mcp} + \texttt{lean\_multi\_attempt}; goal-state
inspection via \texttt{lean\_goal}; tactic-search via
\texttt{lean\_leansearch} / \texttt{lean\_loogle}; substrate-discovery
via \texttt{lean\_local\_search}; type-signature lookup via
\texttt{lean\_hover\_info}; build verification via
\texttt{lean\_verify} / \texttt{lake build}). The disclosure is
positioned as a transparency commitment, not a prohibition: AI tools
are now standard infrastructure for proof development at the project's
scale, and the disclosure protocol asks reviewers to evaluate the
substantive content rather than the tool stack used to author it.

%% ─────────────────────────────────────────────────────────────────
%% §8 Limitations and explicit deferrals
%% ─────────────────────────────────────────────────────────────────

\section{Limitations and explicit deferrals}
\label{sec:limitations}

We list explicitly what the I3 release does \emph{not} carry, so that
downstream consumers and Mathlib4 reviewers can calibrate expectations.

\begin{enumerate}
\item \textbf{No quantitative L${}^2$-isometric construction yet.}
  \texttt{StochasticIntegral.lean} ships a \texttt{Prop := True}
  predicate body. Lifting to the canonical Mathlib-grade L${}^2$-isometric
  construction is the upstream PR target.
\item \textbf{No quantitative Doob--Meyer decomposition yet.}
  \texttt{Semimartingale.lean} ships the predicate but not the
  decomposition existence + uniqueness modulo indistinguishability.
\item \textbf{No quantitative C${}^2$ change-of-variables yet.}
  \texttt{ItoLemma.lean} ships the predicate but not the substantive
  change-of-variables formula.
\item \textbf{No Stratonovich integral.} The Stratonovich integral is
  a related but distinct construction (anti-symmetrization of the
  It\^o integral) and is out of scope for the Phase~6o.$\zeta$
  release.
\item \textbf{No quantitative Cram\'er, Sanov, or Varadhan content
  yet.} Each of the LDP modules ships a predicate-substrate identifier
  with non-vacuity witness; the Cram\'er upper bound, Sanov's
  method-of-types theorem, the contraction principle, the Esscher-tilted
  lower bound, and Varadhan's lemma in their substantive forms are all
  upstream PR targets.
\item \textbf{No infinite-alphabet Sanov.} The substrate-data
  predicate \texttt{IsFiniteAlphabetSanov} is restricted to finite
  alphabets; the infinite-alphabet generalization (Polish-space
  empirical-measure LDP) is a downstream-of-Mathlib-merge target.
\item \textbf{No multi-type LDP.} The Polish-space LDP for
  multi-type empirical measures (e.g.\ Sanov for empirical pair-types
  in Markov-chain LDPs) is out of scope.
\item \textbf{No upstream Mathlib4 PRs yet.} Per
  \S\ref{sec:mathlib-upstream}, the upstream cycle is staged for
  the post-Phase~6o window. The library is currently
  project-internal infrastructure.
\item \textbf{Phase~6o Wave~4a is async-blocked.} The Phase~6o
  Wave~4a substrate-derived $\Lambda_J$ vs $\Lambda_{\rm HK}$ refactor
  --- which is a natural consumer of the I3 LDP substrate via its
  cross-bridge into D5 §11 --- is currently blocked on a
  primary-source dispatch on the Finazzi--Liberati--Sindoni
  acoustic-BEC depletion-factor literature
  (\texttt{Lit-Search/Tasks/submitted/20260508\_FLS\_BEC\_depletion\_factor\_lambda\_substrate.md}).
  The asynchronous return will determine whether the I3 §5
  \texttt{LDPCompatibleSKEFT} typeclass receives a second concrete
  substrate witness on the FLS BEC side or remains at the single
  centered-Gaussian witness.
\end{enumerate}

\subsection{Explicit non-claims}

For the avoidance of doubt: the I3 release does \emph{not} claim to
formalize the It\^o stochastic integral, It\^o's lemma, Novikov's
condition, the Cram\'er upper or lower bound, Sanov's theorem, the
contraction principle, or Varadhan's lemma at the substantive
quantitative level. The release claims to formalize the
\emph{predicate-substrate scaffolding} for these objects, plus a
substantive cross-bridge typeclass that connects the scaffolding to
the program's existing Glorioso--Liu axiomatic content. The
substantive quantitative content is the upstream-PR target, and the
substrate-scaffolding-vs-quantitative-content separation is the
release's explicit engineering choice.

%% ─────────────────────────────────────────────────────────────────
%% §9 Reproducibility and availability
%% ─────────────────────────────────────────────────────────────────

\section{Reproducibility and availability}
\label{sec:availability}

\subsection{License and source}

The library is released under the Apache~2.0 license alongside the
SK--EFT Hawking project's main codebase. The 12 modules live at
\texttt{lean/SKEFTHawking/Itô/} and \texttt{lean/SKEFTHawking/LDP/}
in the project repository.

\subsection{Build}

The library builds against Lean~4 v4.29.0 (per the project's
\texttt{lean-toolchain}) and Mathlib4 pinned at the project's
\texttt{lakefile.toml} commit. From the \texttt{lean/} subdirectory:

\begin{lstlisting}[language=]
lake build SKEFTHawking.LDP.LDPCompatibleSKEFT
\end{lstlisting}

builds the \texttt{LDPCompatibleSKEFT} module and its 11 dependencies
in one command. The full project builds via \texttt{lake build} (no
target) at $\sim$8500 jobs on the project's current state.

\subsection{Verification}

Each headline theorem (the wave-closure summaries; the
\texttt{LDPCompatibleSKEFT} concrete instance) verifies under the
standard Lean~4 kernel plus
\texttt{[propext, Classical.choice, Quot.sound]} via
\texttt{lean\_verify}. Zero \texttt{sorry}; zero new axioms beyond
the standard kernel basis.

\subsection{Contributing}

Pre-Mathlib-upstream-cycle, contributions to the I3 module set
proceed through the SK--EFT Hawking project's normal contribution
channels (per \texttt{CLAUDE.md} pipeline-invariant discipline + the
project's wave-execution pipeline). Post-Mathlib-upstream-cycle, each
upstream-PR-merged module's project-internal copy will be retired and
the canonical content lifted from the upstream Mathlib4 location.

\subsection{Replication}

A reader who wishes to re-verify the library's claims independently
should:
\begin{enumerate}
\item Clone the project repository.
\item Install \texttt{elan} (the Lean~4 version manager) and run
  \texttt{lake update} from the \texttt{lean/} subdirectory.
\item Run \texttt{lake build SKEFTHawking.Itô.Novikov} +
  \texttt{lake build SKEFTHawking.LDP.LDPCompatibleSKEFT}; both
  commands should complete without errors.
\item Optionally run the project's \texttt{lean\_verify} pass on
  the wave-closure summary theorems to confirm standard-kernel-only
  axiom dependence.
\end{enumerate}

\subsection{Citation request}

Authors who use the I3 library in subsequent work are asked to cite
this paper plus the SK--EFT Hawking project's main strategy
document~\cite{Roehm2026Strategy}. Authors who lift the I3 module
content into upstream Mathlib4 PRs are asked to retain the
``Phase~6o.$\zeta$, SK--EFT Hawking project'' attribution in the
PR commit message and to coordinate via the project's contribution
channels for the lift cadence.

%% ─────────────────────────────────────────────────────────────────
%% §10 Conclusion
%% ─────────────────────────────────────────────────────────────────

\section{Conclusion}
\label{sec:conclusion}

We have described a 12-module formally verified Lean~4 library that
establishes predicate-level scaffolding for two foundational pieces
of probability theory currently absent from Mathlib4: the It\^o
stochastic integral and the large-deviation principle. The library
ships under Apache~2.0; verifies standard-kernel-only via
\texttt{lean\_verify}; carries zero \texttt{sorry} and zero new
axioms; couples to the SK--EFT Hawking program's existing
Glorioso--Liu monotonicity axiomatic via the substantive
\texttt{LDPCompatibleSKEFT} typeclass cross-bridge; and is
positioned for a $\sim$11--13-PR upstream Mathlib4 cycle staged for
the post-Phase~6o window. The substrate-vs-quantitative-content engineering choice that drives
the release scope (substrate-data predicates ship now; substantive
quantitative content lifts in the upstream PR cycle) we treat as
continuous with the staged-substrate practice already common in the
Mathlib4 community: the Doob martingale convergence formalization
of Ying and Degenne (CPP~2023)~\cite{YingDegenne2022Martingales}
shipped predicate scaffolding for filtrations, adapted/predictable
processes, stopping times, and sub/super-martingales as a working
substrate, with substantive convergence content lifted in follow-up
mathlib PRs; the Brownian-motion port of Degenne et al.\
(2025)~\cite{Degenne2025BrownianMotion} likewise staged
Carath\'eodory + Kolmogorov + Kolmogorov--Chentsov + Gaussian-measure
substrates before defining Brownian motion proper. I3 follows the
same pattern in a new sub-domain (the It\^o stochastic integral
and the large-deviation principle) and we document the pattern
explicitly here so that the program's Tier~3 software-paper releases
make the workflow more legible for future formal-verification
software releases.

\begin{thebibliography}{99}

\bibitem{Roehm2026Strategy} J.~Roehm,
\emph{SK--EFT Hawking 14-bundle publication architecture (project paper-strategy frame)},
SK--EFT Hawking project document
\texttt{docs/PAPER\_STRATEGY.md}, in preparation (2026).

\bibitem{KaratzasShreve1991BMSC} I.~Karatzas and S.~E.~Shreve,
\emph{Brownian Motion and Stochastic Calculus}, 2nd ed.,
Graduate Texts in Mathematics \textbf{113}, Springer (1991).

\bibitem{RevuzYor1999CMSM} D.~Revuz and M.~Yor,
\emph{Continuous Martingales and Brownian Motion}, 3rd ed.,
Grundlehren der mathematischen Wissenschaften \textbf{293}, Springer (1999).

\bibitem{DemboZeitouni1998LDT} A.~Dembo and O.~Zeitouni,
\emph{Large Deviations Techniques and Applications}, 2nd ed.,
Applications of Mathematics \textbf{38}, Springer (1998).

\bibitem{denHollander2000LDPs} F.~den Hollander,
\emph{Large Deviations}, Fields Institute Monographs \textbf{14},
American Mathematical Society (2000).

\bibitem{FalascoEsposito2025RMP} G.~Falasco and M.~Esposito,
\emph{Macroscopic stochastic thermodynamics},
Reviews of Modern Physics \textbf{97}, 015002 (2025).
DOI: \href{https://doi.org/10.1103/RevModPhys.97.015002}{10.1103/RevModPhys.97.015002}.

\bibitem{Degenne2025MarkovKernels} R.~Degenne,
\emph{Markov kernels in Mathlib's probability library},
arXiv:2510.04070 (2025).
DOI: \href{https://arxiv.org/abs/2510.04070}{arXiv:2510.04070}.
(Primary-source-verified Phase 7 absorption Session 5 2026-05-08
via direct arXiv abstract fetch; supersedes the project's
carry-forward Phase-6n-era ``Marion 2025 MarkovKernels''
misattribution.)

\bibitem{Degenne2025BrownianMotion} R.~Degenne, D.~Ledvinka,
E.~Marion, and P.~Pfaffelhuber,
\emph{Formalization of Brownian motion in Lean},
arXiv:2511.20118 (2025).
DOI: \href{https://arxiv.org/abs/2511.20118}{arXiv:2511.20118}.

\bibitem{mathlib4_2020} The Mathlib Community,
``The Lean Mathematical Library'',
\emph{Proceedings of the 9th ACM SIGPLAN International Conference
on Certified Programs and Proofs (CPP 2020)},
ACM, 367--381 (2020).
DOI: \href{https://doi.org/10.1145/3372885.3373824}{10.1145/3372885.3373824}.

\bibitem{YingDegenne2022Martingales} K.~Ying and R.~Degenne,
``A Formalization of Doob's Martingale Convergence Theorems in mathlib'',
\emph{Proceedings of the 12th ACM SIGPLAN International Conference on
Certified Programs and Proofs (CPP 2023)},
ACM, 334--347 (2023).
arXiv:2212.05578.
DOI: \href{https://doi.org/10.1145/3573105.3575675}{10.1145/3573105.3575675}.

\end{thebibliography}

\end{document}
